\section{动态Prompt输出头容量对照}
\label{sec:output-head-appendix}

输出头压缩实验保留$D=768$的Workspace，并将动态Prompt输出头改为秩256的低秩映射。表\ref{tab:output-head}显示，该改动减少21.84\%的可训练参数，同时使Overall下降3.36个百分点，Free-form下降4.76个百分点。该压缩配置未达到预期的精度与参数量折中，因此主实验采用默认输出头。该对照考察的是固定适配宽度下的一种输出瓶颈设置。

\begin{table}[t]
    \centering
    \caption{PathVQA Validation上的动态Prompt输出头容量对照（seed 44，\%）。}
    \label{tab:output-head}
    \small
    \begin{tabular}{@{}lrrrr@{}}
        \toprule
        配置 & 可训练参数 & Overall & Yes/No & Free-form \\
        \midrule
        QDPT（默认输出头） & 7,805,184 & 60.78 & 92.74 & 28.91 \\
        QDPT（秩256输出头） & 6,100,736 & 57.42 & 90.78 & 24.15 \\
        \bottomrule
    \end{tabular}
\end{table}
